\subsection*{5.2 Double Integrals}
\begin{conceptbox}
    A double integral over a region $R$ is:
    \begin{equation*}
        \iint_R f(x,y)\, dA
    \end{equation*}
\end{conceptbox}

This represents accumulation over an area.
\subsubsection*{Iterated Integrals}

Double integrals are computed as iterated integrals:
\begin{equation*}
    \int_a^b \int_c^d f(x,y)\, dx\, dy
\end{equation*}

\begin{examplebox}
    Evaluate the accumulation of the function $f(x,y)=x+y$ over the region $R=[0,1]\times[0,1]$.

    \textbf{Step 1: Conceptual understanding (before writing the integral)}

    We interpret this as adding up the values of $x+y$ over the entire square region.

    Imagine dividing the region into many tiny rectangles of area $\Delta A$.

    The total accumulation is approximated by:
    \begin{equation*}
        \sum f(x_i,y_i)\,\Delta A
    \end{equation*}

    As the rectangles become infinitely small, this becomes a double integral.

    \textbf{Step 2: Write the double integral}

    \begin{align*}
        \iint_R (x+y)\, dA
    \end{align*}

    Since $R$ is a rectangle, we convert it into an iterated integral:

    \begin{align*}
        \int_0^1 \int_0^1 (x+y)\, dx\, dy
    \end{align*}

    \textbf{Step 3: Integrate with respect to $x$}

    Treat $y$ as constant:
    \begin{align*}
        \int_0^1 (x+y)\, dx
         & = \int_0^1 x\,dx + \int_0^1 y\,dx           \\
         & = \left[\frac{x^2}{2}\right]_0^1 + y[x]_0^1 \\
         & = \frac{1}{2} + y
    \end{align*}

    \textbf{Step 4: Integrate with respect to $y$}

    \begin{align*}
        \int_0^1 \left(\frac{1}{2} + y\right) dy
         & = \left[\frac{y}{2} + \frac{y^2}{2}\right]_0^1 \\
         & = \frac{1}{2} + \frac{1}{2}                    \\
         & = 1
    \end{align*}

    \textbf{Final Answer:} $\boxed{1}$
\end{examplebox}

\begin{examplebox}
    We are given the function:
    \begin{equation*}
        f(x,y) = 1
    \end{equation*}
    over the region:
    \begin{equation*}
        R = [0,2] \times [0,3]
    \end{equation*}

    This region represents a rectangle in the $xy$-plane with:
    \begin{itemize}
        \item width in the $x$-direction: $2 - 0 = 2$
        \item height in the $y$-direction: $3 - 0 = 3$
    \end{itemize}

    Thus, geometrically, we already expect the area to be:
    \[
        \text{Area} = 2 \times 3 = 6
    \]

    \textbf{Step 1: Interpret the double integral}

    The double integral
    \begin{equation*}
        \iint_R 1\, dA
    \end{equation*}
    represents the accumulation of the constant value $1$ over the region $R$.

    Conceptually, this means:
    \begin{itemize}
        \item we are summing many small pieces of area $\Delta A$
        \item each piece contributes exactly its area (since $f=1$)
    \end{itemize}

    So this integral directly computes the \textbf{area of the region}.

    \textbf{Step 2: Convert to an iterated integral}

    Since $R$ is a rectangle, we write:
    \begin{align*}
        \iint_R 1\, dA = \int_0^3 \int_0^2 1\, dx\, dy
    \end{align*}

    \textbf{Step 3: Integrate with respect to $x$}

    Treat $y$ as constant:
    \begin{align*}
        \int_0^2 1\, dx = [x]_0^2 = 2
    \end{align*}

    \textbf{Step 4: Integrate with respect to $y$}

    \begin{align*}
        \int_0^3 2\, dy = [2y]_0^3 = 6
    \end{align*}

    \textbf{Final Answer:}
    \begin{equation*}
        \boxed{6}
    \end{equation*}
\end{examplebox}